\documentclass[11pt,a4paper,modern]{FFExercises}
\usepackage[english,greek]{babel}
\usepackage[utf8]{inputenc}
\usepackage{nimbusserif}
\usepackage[T1]{fontenc}
\usepackage{amsmath}
\let\myBbbk\Bbbk
\let\Bbbk\relax
\usepackage[amsbb,subscriptcorrection,zswash,mtpcal,mtphrb,mtpfrak]{mtpro2}
\usepackage{graphicx,multicol,multirow,enumitem,tabularx,mathimatika,gensymb,venndiagram,hhline,longtable,tkz-euclide,fontawesome5,eurosym,tcolorbox,tabularray,tikzpagenodes,relsize,siunitx,eurosym}
\definecolor{xrwma}{HTML}{aa1212}
\usetikzlibrary{calc}
\usetikzlibrary{positioning}
\tcbuselibrary{skins,theorems,breakable}
\renewcommand{\textstigma}{\textsigma\texttau}
\renewcommand{\textdexiakeraia}{}
\usepackage{svg}
\ekthetesdeiktes
\begin{document}

\titlos{Άλγεβρα}{Α' Λυκείου}{Γεωμετρική Πρόοδος -- Λύσεις}

%%%%%%%%%%%%%%%%%%%%%%%%%%%%%%%%%%%%%%%%%%%%%%%%%%%%%%%%%%%%%%%%%%%%%%%%%%%%
% ΟΜΑΔΑ Α: Βασικοί ορισμοί και αναγνώριση Γ.Π.
%%%%%%%%%%%%%%%%%%%%%%%%%%%%%%%%%%%%%%%%%%%%%%%%%%%%%%%%%%%%%%%%%%%%%%%%%%%%

\paragraph{Ομάδα Α: Βασικοί ορισμοί και αναγνώριση Γ.Π.}

\askhsh
\textbf{Λύση}

Ελέγχουμε αν ο λόγος διαδοχικών όρων είναι σταθερός:
\begin{alist}
  \item $\dfrac{6}{2} = 3$, $\dfrac{18}{6} = 3$, $\dfrac{54}{18} = 3$. \textbf{Είναι Γ.Π.} με $\lambda = 3$.
  \item $\dfrac{10}{5} = 2$, $\dfrac{20}{10} = 2$, $\dfrac{40}{20} = 2$. \textbf{Είναι Γ.Π.} με $\lambda = 2$.
  \item $\dfrac{-2}{1} = -2$, $\dfrac{4}{-2} = -2$, $\dfrac{-8}{4} = -2$. \textbf{Είναι Γ.Π.} με $\lambda = -2$.
  \item $\dfrac{6}{3} = 2$, $\dfrac{12}{6} = 2$, $\dfrac{15}{12} = 1{,}25 \neq 2$. \textbf{Δεν είναι Γ.Π.}
  \item $\dfrac{27}{81} = \dfrac{1}{3}$, $\dfrac{9}{27} = \dfrac{1}{3}$, $\dfrac{3}{9} = \dfrac{1}{3}$. \textbf{Είναι Γ.Π.} με $\lambda = \dfrac{1}{3}$.
  \item $\dfrac{1/4}{1/2} = \dfrac{1}{2}$, $\dfrac{1/8}{1/4} = \dfrac{1}{2}$, $\dfrac{1/16}{1/8} = \dfrac{1}{2}$. \textbf{Είναι Γ.Π.} με $\lambda = \dfrac{1}{2}$.
\end{alist}

\askhsh
\textbf{Λύση}

Ο λόγος Γ.Π. είναι $\lambda = \dfrac{\alpha_2}{\alpha_1}$:
\begin{alist}
  \item $\lambda = \dfrac{12}{4} = 3$
  \item $\lambda = \dfrac{50}{100} = \dfrac{1}{2}$
  \item $\lambda = \dfrac{6}{-2} = -3$
  \item $\lambda = \dfrac{1}{1/3} = 3$
\end{alist}

\askhsh
\textbf{Λύση}
\begin{alist}
  \item Με $\alpha_1 = 3$ και $\lambda = 2$:
  \[
  \alpha_1 = 3,\ \alpha_2 = 6,\ \alpha_3 = 12,\ \alpha_4 = 24,\ \alpha_5 = 48
  \]
  
  \item Λύνουμε $\alpha_\nu = 192$:
  \[
  3 \cdot 2^{\nu-1} = 192 \Rightarrow 2^{\nu-1} = 64 = 2^6 \Rightarrow \nu - 1 = 6 \Rightarrow \nu = 7
  \]
\end{alist}

\askhsh
\textbf{Λύση}
\begin{alist}
  \item Με $\alpha_1 = 64$ και $\lambda = \dfrac{1}{2}$:
  \[
  \alpha_1 = 64,\ \alpha_2 = 32,\ \alpha_3 = 16,\ \alpha_4 = 8,\ \alpha_5 = 4,\ \alpha_6 = 2
  \]
  
  \item Λύνουμε $\alpha_\nu > 1$:
  \[
  64 \cdot \left(\frac{1}{2}\right)^{\nu-1} > 1 \Rightarrow \left(\frac{1}{2}\right)^{\nu-1} > \frac{1}{64}
  \]
  \[
  \left(\frac{1}{2}\right)^{\nu-1} > \left(\frac{1}{2}\right)^6 \Rightarrow \nu - 1 < 6 \Rightarrow \nu < 7
  \]
  Άρα $\nu \in \{1, 2, 3, 4, 5, 6\}$.
\end{alist}

\askhsh
\textbf{Λύση}
\begin{alist}
  \item Από $\alpha_2 = 6$ και $\alpha_3 = 18$:
  \[
  \lambda = \frac{\alpha_3}{\alpha_2} = \frac{18}{6} = 3, \quad \alpha_1 = \frac{\alpha_2}{\lambda} = \frac{6}{3} = 2
  \]
  
  \item Έχουμε $\alpha_\nu = 2 \cdot 3^{\nu-1}$, άρα:
  \[
  \beta_\nu = \log_3(2 \cdot 3^{\nu-1}) = \log_3 2 + (\nu-1)
  \]
  Η ακολουθία $(\beta_\nu)$ είναι Α.Π. με $\beta_1 = \log_3 2$ και διαφορά $\omega = 1$.
\end{alist}

\askhsh
\textbf{Λύση}
\begin{alist}
  \item Από $\alpha_6 = \alpha_4 \cdot \lambda^2$:
  \[
  96 = 24 \cdot \lambda^2 \Rightarrow \lambda^2 = 4 \Rightarrow \lambda = \pm 2
  \]
  
  \item Για $\lambda = 2$: $\alpha_4 = \alpha_1 \cdot 2^3 \Rightarrow 24 = 8\alpha_1 \Rightarrow \alpha_1 = 3$
  
  \item Λύνουμε $\alpha_\nu < 1000$ με $\alpha_1 = 3$, $\lambda = 2$:
  \[
  3 \cdot 2^{\nu-1} < 1000 \Rightarrow 2^{\nu-1} < 333{,}3
  \]
  Αφού $2^8 = 256 < 333{,}3 < 512 = 2^9$: $\nu - 1 \leq 8 \Rightarrow \nu \leq 9$.
  
  Άρα $\nu \in \{1, 2, \ldots, 9\}$.
\end{alist}

%%%%%%%%%%%%%%%%%%%%%%%%%%%%%%%%%%%%%%%%%%%%%%%%%%%%%%%%%%%%%%%%%%%%%%%%%%%%
% ΟΜΑΔΑ Β: Υπολογισμός ν-οστού όρου
%%%%%%%%%%%%%%%%%%%%%%%%%%%%%%%%%%%%%%%%%%%%%%%%%%%%%%%%%%%%%%%%%%%%%%%%%%%%

\paragraph{Ομάδα Β: Ο ν-οστός όρος Γ.Π.}

\askhsh
\textbf{Λύση}

Χρησιμοποιούμε τον τύπο $\alpha_\nu = \alpha_1 \cdot \lambda^{\nu-1}$:
\begin{alist}
  \item $\alpha_{10} = 2 \cdot 2^9 = 2^{10} = 1024$
  \item $\alpha_{10} = 3 \cdot 3^9 = 3^{10} = 59049$
  \item $\alpha_{10} = 1 \cdot (-2)^9 = -512$
  \item $\alpha_{10} = 1024 \cdot \left(\dfrac{1}{2}\right)^9 = 1024 \cdot \dfrac{1}{512} = 2$
\end{alist}

\askhsh
\textbf{Λύση}
\begin{alist}
  \item $\alpha_\nu = \alpha_1 \cdot \lambda^{\nu-1} = 5 \cdot 3^{\nu-1}$
  \item $\alpha_7 = 5 \cdot 3^6 = 5 \cdot 729 = 3645$
\end{alist}

\askhsh
\textbf{Λύση}
\begin{alist}
  \item $\alpha_\nu = 2 \cdot (-3)^{\nu-1}$
  \item $\alpha_8 = 2 \cdot (-3)^7 = 2 \cdot (-2187) = -4374$
  \item Για $\nu = 100$: $(-3)^{99} = (-1)^{99} \cdot 3^{99} = -3^{99} < 0$.
  
  Άρα $\alpha_{100} = 2 \cdot (-3)^{99} < 0$.
\end{alist}

\askhsh
\textbf{Λύση}

Από $\alpha_5 = \alpha_1 \cdot \lambda^4$:
\[
48 = \alpha_1 \cdot 2^4 = 16 \alpha_1 \Rightarrow \alpha_1 = 3
\]
\[
\alpha_{10} = \alpha_1 \cdot \lambda^9 = 3 \cdot 2^9 = 3 \cdot 512 = 1536
\]

\askhsh
\textbf{Λύση}
\begin{alist}
  \item Από $\alpha_6 = \alpha_3 \cdot \lambda^3$:
  \[
  96 = 12 \cdot \lambda^3 \Rightarrow \lambda^3 = 8 \Rightarrow \lambda = 2
  \]
  
  \item Από $\alpha_3 = \alpha_1 \cdot \lambda^2$:
  \[
  12 = \alpha_1 \cdot 4 \Rightarrow \alpha_1 = 3
  \]
  
  \item $\alpha_{10} = 3 \cdot 2^9 = 3 \cdot 512 = 1536$
\end{alist}

\askhsh
\textbf{Λύση}
\begin{alist}
  \item Από $\alpha_5 = \alpha_2 \cdot \lambda^3$:
  \[
  80 = 10 \cdot \lambda^3 \Rightarrow \lambda^3 = 8 \Rightarrow \lambda = 2
  \]
  
  \item $\alpha_8 = \alpha_5 \cdot \lambda^3 = 80 \cdot 8 = 640$
  
  \item $\alpha_1 = \dfrac{\alpha_2}{\lambda} = \dfrac{10}{2} = 5$
  
  $S_6 = 5 \cdot \dfrac{2^6 - 1}{2 - 1} = 5 \cdot 63 = 315$
\end{alist}

\askhsh
\textbf{Λύση}
\begin{alist}
  \item Σε Γ.Π. ισχύει $\alpha_1 \cdot \alpha_5 = \alpha_3^2 \Rightarrow 81 = \alpha_3^2 \Rightarrow \alpha_3 = 9$
  
  \item Με $\alpha_1 + \alpha_5 = 30$ και $\alpha_1 \cdot \alpha_5 = 81$, οι $\alpha_1$, $\alpha_5$ είναι ρίζες της:
  \[
  t^2 - 30t + 81 = 0 \Rightarrow t = \frac{30 \pm \sqrt{900-324}}{2} = \frac{30 \pm 24}{2}
  \]
  Άρα $\alpha_1 = 3$, $\alpha_5 = 27$ ή αντίστροφα.
\end{alist}

\askhsh
\textbf{Λύση}
\begin{alist}
  \item Σε Γ.Π. ισχύει $\alpha_3 \cdot \alpha_7 = \alpha_5^2 \Rightarrow 144 = \alpha_5^2 \Rightarrow \alpha_5 = 12$ (αφού $\alpha_5 > 0$)
  
  \item Με $\lambda = 2$: $\alpha_5 = \alpha_1 \cdot 2^4 \Rightarrow 12 = 16\alpha_1 \Rightarrow \alpha_1 = \dfrac{3}{4}$
\end{alist}

%%%%%%%%%%%%%%%%%%%%%%%%%%%%%%%%%%%%%%%%%%%%%%%%%%%%%%%%%%%%%%%%%%%%%%%%%%%%
% ΟΜΑΔΑ Γ: Άθροισμα όρων Γ.Π.
%%%%%%%%%%%%%%%%%%%%%%%%%%%%%%%%%%%%%%%%%%%%%%%%%%%%%%%%%%%%%%%%%%%%%%%%%%%%

\paragraph{Ομάδα Γ: Άθροισμα όρων Γ.Π.}

\askhsh
\textbf{Λύση}

Χρησιμοποιούμε $S_\nu = \alpha_1 \cdot \dfrac{\lambda^\nu - 1}{\lambda - 1}$:
\begin{alist}
  \item $\alpha_1 = 1$, $\lambda = 2$, $\nu = 10$:
  \[
  S_{10} = 1 \cdot \frac{2^{10} - 1}{2 - 1} = 1024 - 1 = 1023
  \]
  
  \item $\alpha_1 = 3$, $\lambda = 3$, $\nu = 8$:
  \[
  S_8 = 3 \cdot \frac{3^8 - 1}{3 - 1} = 3 \cdot \frac{6561 - 1}{2} = 3 \cdot 3280 = 9840
  \]
  
  \item $\alpha_1 = 1$, $\lambda = -2$, $\alpha_\nu = 512 = (-2)^9$, άρα $\nu = 10$:
  \[
  S_{10} = 1 \cdot \frac{(-2)^{10} - 1}{-2 - 1} = \frac{1024 - 1}{-3} = -341
  \]
  
  \item $\alpha_1 = 1$, $\lambda = \dfrac{1}{2}$, $\alpha_\nu = \dfrac{1}{256} = \left(\dfrac{1}{2}\right)^8$, άρα $\nu = 9$:
  \[
  S_9 = 1 \cdot \frac{1 - (1/2)^9}{1 - 1/2} = 2 \cdot \frac{511}{512} = \frac{511}{256}
  \]
\end{alist}

\askhsh
\textbf{Λύση}
\begin{alist}
  \item $S_6 = 4 \cdot \dfrac{3^6 - 1}{3 - 1} = 4 \cdot \dfrac{728}{2} = 1456$
  
  \item Λύνουμε $4 \cdot \dfrac{3^\nu - 1}{2} = 1456$:
  \[
  3^\nu - 1 = 728 \Rightarrow 3^\nu = 729 = 3^6 \Rightarrow \nu = 6
  \]
\end{alist}

\askhsh
\textbf{Λύση}
\begin{alist}
  \item $S_8 = 2 \cdot \dfrac{2^8 - 1}{2 - 1} = 2 \cdot 255 = 510$
  
  \item Λύνουμε $S_\nu = 2046$:
  \[
  2 \cdot \frac{2^\nu - 1}{1} = 2046 \Rightarrow 2^\nu - 1 = 1023 \Rightarrow 2^\nu = 1024 = 2^{10} \Rightarrow \nu = 10
  \]
\end{alist}

\askhsh
\textbf{Λύση}
\begin{alist}
  \item $S_5 = 6 \cdot \dfrac{1 - (1/3)^5}{1 - 1/3} = 6 \cdot \dfrac{3}{2} \cdot \dfrac{242}{243} = \dfrac{242}{27}$
  
  \item Λύνουμε $\alpha_\nu > \dfrac{1}{100}$:
  \[
  6 \cdot \left(\frac{1}{3}\right)^{\nu-1} > \frac{1}{100} \Rightarrow \left(\frac{1}{3}\right)^{\nu-1} > \frac{1}{600}
  \]
  Αφού $3^5 = 243 < 600 < 729 = 3^6$: $\nu - 1 < 6 \Rightarrow \nu < 7$.
  
  Άρα $\nu \in \{1, 2, 3, 4, 5, 6\}$.
\end{alist}

\askhsh
\textbf{Λύση}
\begin{alist}
  \item $S_{10} = 1 \cdot \dfrac{1 - (-2)^{10}}{1 - (-2)} = \dfrac{1 - 1024}{3} = -341$
  
  \item $|\alpha_\nu| = |2 \cdot (-3)^{\nu-1}| = 2 \cdot 2^{\nu-1}$.
  
  Έχουμε $\dfrac{|\alpha_{\nu+1}|}{|\alpha_\nu|} = \dfrac{2 \cdot 2^\nu}{2 \cdot 2^{\nu-1}} = 2$.
  
  Άρα η $(|\alpha_\nu|)$ είναι Γ.Π. με λόγο $\lambda = 2$.
\end{alist}

\askhsh
\textbf{Λύση}
\begin{alist}
  \item $\alpha_4 = \alpha_1 \cdot \lambda^3 \Rightarrow 54 = 2 \cdot \lambda^3 \Rightarrow \lambda = 3$
  
  \item $S_6 = 2 \cdot \dfrac{3^6 - 1}{3 - 1} = 2 \cdot \dfrac{728}{2} = 728$
  
  \item $\alpha_6 = 2 \cdot 3^5 = 486$, άρα $2\alpha_6 = 972$.
  
  Συγκρίνουμε: $S_6 = 728 < 972 = 2\alpha_6$.
\end{alist}

\askhsh
\textbf{Λύση}
\begin{alist}
  \item $S_4 = \alpha_1 \cdot \dfrac{3^4 - 1}{3 - 1} = 40\alpha_1 = 40 \Rightarrow \alpha_1 = 1$
  
  \item $S_6 - S_4 = \alpha_5 + \alpha_6 = 1 \cdot 3^4 + 1 \cdot 3^5 = 81 + 243 = 324$
\end{alist}

\askhsh
\textbf{Λύση}

Από $\alpha_\nu = \alpha_1 \cdot \lambda^{\nu-1}$:
\[
256 = 1 \cdot \lambda^{\nu-1} \Rightarrow \lambda^{\nu-1} = 256
\]

Από $S_\nu = \dfrac{\lambda^\nu - 1}{\lambda - 1}$:
\[
511 = \frac{\lambda^\nu - 1}{\lambda - 1}
\]

Αφού $\lambda^{\nu-1} = 256$, τότε $\lambda^\nu = 256\lambda$.
\[
511(\lambda - 1) = 256\lambda - 1 \Rightarrow 255\lambda = 510 \Rightarrow \lambda = 2
\]

Από $2^{\nu-1} = 256 = 2^8$: $\nu = 9$.

%%%%%%%%%%%%%%%%%%%%%%%%%%%%%%%%%%%%%%%%%%%%%%%%%%%%%%%%%%%%%%%%%%%%%%%%%%%%
% ΟΜΑΔΑ Δ: Εύρεση παραμέτρου
%%%%%%%%%%%%%%%%%%%%%%%%%%%%%%%%%%%%%%%%%%%%%%%%%%%%%%%%%%%%%%%%%%%%%%%%%%%%

\paragraph{Ομάδα Δ: Εύρεση παραμέτρου}

\askhsh
\textbf{Λύση}
\begin{alist}
  \item $x^2 = 2 \cdot 8 = 16 \Rightarrow x = 4$ (αφού $x > 0$)
  
  \item $\lambda = \dfrac{x}{2} = \dfrac{4}{2} = 2$
  
  \item Για Α.Π.: $x = \dfrac{2 + 8}{2} = 5 \neq 4$. Άρα \textbf{όχι}.
\end{alist}

\askhsh
\textbf{Λύση}
\begin{alist}
  \item $x^2 = 3 \cdot 12 = 36 \Rightarrow x = 6$ (αφού $x > 0$)
  
  \item $\lambda = \dfrac{6}{3} = 2$.
  
  $\alpha_5 = 3 \cdot 2^4 = 48$
\end{alist}

\askhsh
\textbf{Λύση}
\begin{alist}
  \item $(x+1)^2 = (x-1)(x+5) \Rightarrow x^2 + 2x + 1 = x^2 + 4x - 5 \Rightarrow x = 3$
  
  \item Οι τρεις όροι: $2,\ 4,\ 8$ με $\lambda = 2$.
  
  \item $S_5 = 2 \cdot \dfrac{2^5 - 1}{2 - 1} = 2 \cdot 31 = 62$
\end{alist}

\askhsh
\textbf{Λύση}

$\alpha \cdot \beta = 36$ και $\alpha + \beta = 15$.

$t^2 - 15t + 36 = 0 \Rightarrow t = 3$ ή $t = 12$.

Επειδή $\alpha < \beta$: $\alpha = 3$ και $\beta = 12$.

\askhsh
\textbf{Λύση}

Έστω $\dfrac{\alpha}{\lambda},\ \alpha,\ \alpha\lambda$.

Γινόμενο: $\alpha^3 = 216 \Rightarrow \alpha = 6$

Άθροισμα: $\dfrac{6}{\lambda} + 6 + 6\lambda = 21 \Rightarrow 2\lambda^2 - 5\lambda + 2 = 0 \Rightarrow \lambda = 2$ ή $\lambda = \dfrac{1}{2}$

Οι αριθμοί: $3,\ 6,\ 12$.

\askhsh
\textbf{Λύση}

Έστω $\dfrac{\alpha}{\lambda},\ \alpha,\ \alpha\lambda$.

$\alpha(s + 1) = 26$ και $\alpha(s - 1) = 14$ όπου $s = \lambda + \dfrac{1}{\lambda}$.

$2\alpha = 12 \Rightarrow \alpha = 6$, $s = \dfrac{10}{3}$, $\lambda = 3$ ή $\lambda = \dfrac{1}{3}$.

Οι αριθμοί: $2,\ 6,\ 18$.

\askhsh
\textbf{Λύση}

Έστω $\dfrac{\alpha}{\lambda},\ \alpha,\ \alpha\lambda$.

$\alpha^2 = 16 \Rightarrow \alpha = 4$, $\lambda = 2$ ή $\lambda = \dfrac{1}{2}$.

Οι αριθμοί: $2,\ 4,\ 8$.

%%%%%%%%%%%%%%%%%%%%%%%%%%%%%%%%%%%%%%%%%%%%%%%%%%%%%%%%%%%%%%%%%%%%%%%%%%%%
% ΟΜΑΔΑ Ε: Ιδιότητες Γ.Π.
%%%%%%%%%%%%%%%%%%%%%%%%%%%%%%%%%%%%%%%%%%%%%%%%%%%%%%%%%%%%%%%%%%%%%%%%%%%%

\paragraph{Ομάδα Ε: Ιδιότητες Γ.Π.}

\askhsh
\textbf{Λύση}
\begin{alist}
  \item $\alpha_\mu \cdot \alpha_\nu = \alpha_1^2 \lambda^{\mu+\nu-2} = \alpha_1 \cdot \alpha_{\mu+\nu-1}$ \checkmark
  
  \item $\alpha_\nu^2 = \alpha_1^2 \lambda^{2\nu-2} = \alpha_{\nu-1} \cdot \alpha_{\nu+1}$ \checkmark
\end{alist}

\askhsh
\textbf{Λύση}

$\dfrac{\alpha_{\nu+2}}{\alpha_\nu} = \dfrac{\alpha_1 \lambda^{\nu+1}}{\alpha_1 \lambda^{\nu-1}} = \lambda^2$

\askhsh
\textbf{Λύση}
\begin{alist}
  \item $\dfrac{\alpha_8}{\alpha_5} = \lambda^3 \Rightarrow \dfrac{384}{48} = 8 = \lambda^3 \Rightarrow \lambda = 2$
  
  \item $\alpha_5 \cdot \alpha_8 = \alpha_6 \cdot \alpha_7 = 48 \cdot 384 = 18432$
  
  \item $\alpha_5 + \alpha_6 + \alpha_7 + \alpha_8 = 48(1 + 2 + 4 + 8) = 48 \cdot 15 = 720$
\end{alist}

\askhsh
\textbf{Λύση}
\begin{alist}
  \item $\dfrac{\alpha_7}{\alpha_3} = \lambda^4 \Rightarrow \dfrac{128}{8} = 16 = \lambda^4 \Rightarrow \lambda = 2$
  
  $\alpha_3 = \alpha_1 \cdot 4 = 8 \Rightarrow \alpha_1 = 2$
  
  \item $\alpha_3 \cdot \alpha_7 = \alpha_4 \cdot \alpha_6 = 8 \cdot 128 = 1024$
  
  \item $\alpha_\nu = 512 \Rightarrow 2 \cdot 2^{\nu-1} = 512 \Rightarrow 2^\nu = 512 = 2^9 \Rightarrow \nu = 9$
\end{alist}

\askhsh
\textbf{Λύση}

$\dfrac{\alpha_3 + \alpha_5}{\alpha_2 + \alpha_4} = \lambda = 2$

$\alpha_1 \cdot 2 + \alpha_1 \cdot 8 = 30 \Rightarrow \alpha_1 = 3$

%%%%%%%%%%%%%%%%%%%%%%%%%%%%%%%%%%%%%%%%%%%%%%%%%%%%%%%%%%%%%%%%%%%%%%%%%%%%
% ΟΜΑΔΑ ΣΤ: Εφαρμογές
%%%%%%%%%%%%%%%%%%%%%%%%%%%%%%%%%%%%%%%%%%%%%%%%%%%%%%%%%%%%%%%%%%%%%%%%%%%%

\paragraph{Ομάδα ΣΤ: Εφαρμογές -- Πρακτικά προβλήματα}

\askhsh
\textbf{Λύση}
\begin{alist}
  \item $\alpha_9 = 500 \cdot 2^8 = 128\,000$ βακτήρια.
  
  \item $500 \cdot 2^{\nu-1} > 500\,000 \Rightarrow 2^{\nu-1} > 1000$.
  
  $2^{10} = 1024 > 1000 \Rightarrow \nu \geq 11$. Μετά από 10 ώρες.
\end{alist}

\askhsh
\textbf{Λύση}
\begin{alist}
  \item $\alpha_6 = 20\,000 \cdot 0{,}8^5 \approx 6\,553{,}60$ ευρώ.
  
  \item $0{,}8^{\nu-1} < 0{,}25 \Rightarrow \nu \geq 8$. Μετά από 7 χρόνια.
\end{alist}

\askhsh
\textbf{Λύση}
\begin{alist}
  \item $\alpha_6 = 80 \cdot \left(\dfrac{1}{2}\right)^5 = 2{,}5$ γραμμάρια
  
  \item $80 \cdot \left(\dfrac{1}{2}\right)^{\nu-1} < 1 \Rightarrow 2^{\nu-1} > 80$.
  
  $2^7 = 128 > 80 \Rightarrow \nu \geq 8$. Μετά από 14 χρόνια.
\end{alist}

\askhsh
\textbf{Λύση}
\begin{alist}
  \item $\alpha_5 = 16 \cdot \left(\dfrac{3}{4}\right)^4 = \dfrac{81}{16} = 5{,}0625$ m
  
  \item $\left(\dfrac{3}{4}\right)^{\nu-1} < \dfrac{1}{16} \Rightarrow \nu \geq 11$. Μετά την 10η αναπήδηση.
\end{alist}

\askhsh
\textbf{Λύση}
\begin{alist}
  \item $K_4 = 8\,000 \cdot 1{,}05^4 \approx 9\,724{,}05$ ευρώ
  
  \item $1{,}05^\nu = 2 \Rightarrow \nu = \dfrac{\log 2}{\log 1{,}05} \approx 14{,}2$.
  
  Μετά από 15 χρόνια.
\end{alist}

\askhsh
\textbf{Λύση}
\begin{alist}
  \item $S = 1000 \cdot 1{,}04 \cdot \dfrac{1{,}04^5 - 1}{0{,}04} \approx 5\,632{,}98$ ευρώ
  
  \item Η τελική αξία της $k$-οστής κατάθεσης: $1000 \cdot 1{,}04^{6-k}$.
  
  Αυτές αποτελούν Γ.Π. με $\lambda = \dfrac{1}{1{,}04} \approx 0{,}962$.
\end{alist}

%%%%%%%%%%%%%%%%%%%%%%%%%%%%%%%%%%%%%%%%%%%%%%%%%%%%%%%%%%%%%%%%%%%%%%%%%%%%
% ΟΜΑΔΑ Ζ: Συνδυαστικές ασκήσεις
%%%%%%%%%%%%%%%%%%%%%%%%%%%%%%%%%%%%%%%%%%%%%%%%%%%%%%%%%%%%%%%%%%%%%%%%%%%%

\paragraph{Ομάδα Ζ: Συνδυαστικές ασκήσεις}

\askhsh
\textbf{Λύση}
\begin{alist}
  \item $\lambda^3 = 8 \Rightarrow \lambda = 2$. Άρα $x = 4$ και $y = 8$.
  
  \item $S_4 = 2 + 4 + 8 + 16 = 30$
  
  \item Για Α.Π.: $x - 2 = y - x = 16 - y \Rightarrow x = 6{,}y = 10$.
  
  Αλλά $x = 4$ (από Γ.Π.) $\neq 6$. Άρα \textbf{όχι}.
\end{alist}

\askhsh
\textbf{Λύση}
\begin{alist}
  \item $\lambda^3 = 27 \Rightarrow \lambda = 3$. Άρα $x = 6$ και $y = 18$.
  
  \item $\alpha_6 = 2 \cdot 3^5 = 486$
\end{alist}

\askhsh
\textbf{Λύση}

Α.Π.: $\beta = \dfrac{1 + 5}{2} = 3$.

Επαλήθευση Γ.Π.: $1,\ 3,\ 9$ με $\lambda = 3$. \checkmark

\askhsh
\textbf{Λύση}

Α.Π.: $\alpha,\ \alpha+2,\ \alpha+4$ και 4ος = $\alpha + 9$.

Γ.Π.: $(\alpha+2)^2 = \alpha(\alpha+9) \Rightarrow \alpha = \dfrac{4}{5}$.

\askhsh
\textbf{Λύση}

$\alpha_2^2 - 6\alpha_2 + 1 = 0 \Rightarrow \alpha_2 = 3 \pm 2\sqrt{2}$.

$\alpha_3 = 2\alpha_2 - 2 = 4 + 4\sqrt{2}$ (για $\alpha_2 = 3 + 2\sqrt{2}$).

%%%%%%%%%%%%%%%%%%%%%%%%%%%%%%%%%%%%%%%%%%%%%%%%%%%%%%%%%%%%%%%%%%%%%%%%%%%%
% ΟΜΑΔΑ Η: Σωστό - Λάθος
%%%%%%%%%%%%%%%%%%%%%%%%%%%%%%%%%%%%%%%%%%%%%%%%%%%%%%%%%%%%%%%%%%%%%%%%%%%%

\paragraph{Ομάδα Η: Σωστό -- Λάθος}

\askhsh
\textbf{Λύση}
\begin{alist}
  \item \textbf{Σ}
  \item \textbf{Σ}
  \item \textbf{Σ}
  \item \textbf{Λ}
  \item \textbf{Λ}
\end{alist}

\askhsh
\textbf{Λύση}
\begin{alist}
  \item \textbf{Σ}
  \item \textbf{Σ}
  \item \textbf{Σ}
  \item \textbf{Λ}
  \item \textbf{Λ}
\end{alist}

%%%%%%%%%%%%%%%%%%%%%%%%%%%%%%%%%%%%%%%%%%%%%%%%%%%%%%%%%%%%%%%%%%%%%%%%%%%%
% ΟΜΑΔΑ Θ: Αποδείξεις
%%%%%%%%%%%%%%%%%%%%%%%%%%%%%%%%%%%%%%%%%%%%%%%%%%%%%%%%%%%%%%%%%%%%%%%%%%%%

\paragraph{Ομάδα Θ: Αποδείξεις}

\askhsh
\textbf{Λύση}
\begin{alist}
  \item $\dfrac{\alpha_\nu}{\alpha_\mu} = \lambda^{\nu-\mu} \Rightarrow \alpha_\nu = \alpha_\mu \cdot \lambda^{\nu-\mu}$
  
  \item Με $\alpha_3 = 12$ και $\lambda = 2$:
  
  $\alpha_1 = \alpha_3 \cdot 2^{1-3} = 12 \cdot \dfrac{1}{4} = 3$
  
  $\alpha_7 = \alpha_3 \cdot 2^{7-3} = 12 \cdot 16 = 192$
\end{alist}

\askhsh
\textbf{Λύση}

$S_\nu = \alpha_1(1 + \lambda + \ldots + \lambda^{\nu-1})$.

$S_\nu - \lambda S_\nu = \alpha_1(1 - \lambda^\nu) \Rightarrow S_\nu = \alpha_1 \cdot \dfrac{1 - \lambda^\nu}{1 - \lambda}$

\askhsh
\textbf{Λύση}
\begin{alist}
  \item $\beta = \alpha\lambda$, $\gamma = \alpha\lambda^2 \Rightarrow \alpha\gamma = \alpha^2\lambda^2 = \beta^2$
  
  \item $\beta^2 = 3 \cdot 27 = 81 \Rightarrow \beta = 9$ (αφού $\beta > 0$)
\end{alist}

%%%%%%%%%%%%%%%%%%%%%%%%%%%%%%%%%%%%%%%%%%%%%%%%%%%%%%%%%%%%%%%%%%%%%%%%%%%%
% ΟΜΑΔΑ Ι: Πολλαπλής επιλογής
%%%%%%%%%%%%%%%%%%%%%%%%%%%%%%%%%%%%%%%%%%%%%%%%%%%%%%%%%%%%%%%%%%%%%%%%%%%%

\paragraph{Ομάδα Ι: Πολλαπλής επιλογής}

\askhsh
\textbf{Λύση}

$\alpha_4 = 2 \cdot 3^3 = 54$. \textbf{Απάντηση: (γ)}

\askhsh
\textbf{Λύση}

$\lambda^2 = 9 \Rightarrow \lambda = 3$. \textbf{Απάντηση: (β)}

\askhsh
\textbf{Λύση}

$S_5 = 31$. \textbf{Απάντηση: (β)}

\end{document}
